\ProvidesFile{tikzlibrarycartes.code.tex}[2026/05/01 v1.0.0 A tikz library for geometry computations]

\RequirePackage{tikz}
\RequirePackage{tikz-3dplot}
\RequirePackage{ifthen}
\RequirePackage{amsmath}
\RequirePackage{amssymb}
\RequirePackage{bm}

\usetikzlibrary{math,calc,quotes,mc,vc,conics}

% pgfmath-solution to the error:
% Package PGF Math: Could not parse input 'A' as a floating 
% point number, sorry. The unreadable part was near 'A'..
% https://tex.stackexchange.com/questions/455991/pgfmath-function-for-strings-and-numbers
\pgfkeys{
  /pgf/fpu/handlers/invalid number/.code = {%
    \pgfmathfloatparsenumber{3Y0.0e0]}%
  }
}

\makeatletter

% ===============================================
% Caveats
% ===============================================

% function 'intersect/eqs' conflicts with siunitx package

% 在 TikZ 中，函数本身并没有直接的“返回多个值”的机制，因为 TikZ 是一种基于 TeX 的绘图语言，主要用于描述图形，
% 而不是像编程语言那样处理多返回值逻辑。不过，通过一些技巧，你可以在 TikZ 中模拟类似“返回多个值”的效果。
% https://tex.stackexchange.com/questions/510418/tikz-declare-function-with-multiple-outputs

% tikzset 中不要有空行

% 注意: 计算过程是保留坐标单位(pt)的, 所以存在乘除法单位的问题, 首先数值始终携带单位, 
% 在 calc 运算时有的需要转换为标量; 将坐标转换为 pt 值, 数值可能超出限值, 出现 
% Dimension too large 错误, 在计算长度时及时进行缩小 
% https://tex.stackexchange.com/questions/475556/tikz-why-is-dimension-too-large
% 具体方法是修改默认的 1cm, 如: [scale=1.0,x=0.5cm,y=0.5cm]
% 注意此处的变量不要和 tikzpicture 环境重名, 否则被替换掉
% triangle centers: 
% https://mathworld.wolfram.com/BarycentricCoordinates.html

% Syntax of parabola: y = a * x * x + b * x + c
%   \parabola [options] (a,b,c); 
\newcommand\parabola{} % just for safety
\def\parabola[#1]#2(#3,#4,#5){
  \path[smooth,domain=-2:2,#1,variable=\euclidea@var]
    plot({\euclidea@var},
      {(#3)*\euclidea@var*\euclidea@var+(#4)*\euclidea@var+(#5)})
}

% Syntax of ellipse:
%   \ellipse [options] (a,b); 
% wherein, a,b are major/minor semi axis.
\newcommand\ellipse{} % just for safety
\def\ellipse[#1]#2(#3,#4){
  \path[#1] (0,0) ellipse ({#3} and {#4})
}

% Syntax of hyperbola:
%   \hyperbola [options] (a,b); 
% wherein, a,b are major/minor semi axis.
\newcommand\hyperbola{} % just for safety
\def\hyperbola[#1]#2(#3,#4){
  \path[smooth,domain=-1.5:1.5,#1,variable=\euclidea@var]
    plot({-(#3)*cosh(\euclidea@var)},
      {(#4)*sinh(\euclidea@var)})% right arm
    plot({ (#3)*cosh(\euclidea@var)},
      {(#4)*sinh(\euclidea@var)})% left arm
}

% Syntax of asymptote:
%   \asymptote [options] (a,b); 
% wherein, a,b are major/minor semi axis.
\newcommand\asymptote{} % just for safety
\def\asymptote[#1]#2(#3,#4){
  \path[domain=-2:2,#1,variable=\euclidea@var,]
    plot({\euclidea@var},{ (#4)/(#3)*(\euclidea@var)})  
    plot({\euclidea@var},{-(#4)/(#3)*(\euclidea@var)}) 
}

% Syntax: \axes[grid] (xmin:xmax, ymin:ymax)
\newcommand{\axes}[2][]{%
  \axesinternal[#1]#2%
}

% 辅助宏：接收 [options] 和 (content)
\def\axesinternal[#1](#2:#3,#4:#5){%
  % Check if the "grid" option is provided
  \ifthenelse{\equal{#1}{grid}}{%
    \draw[help lines] (#2,#4) grid[step=1] (#3,#5);
  }{}%

  % Axes
  \draw[-latex] ({(#2)-0.5},0) -- ({(#3)+0.5},0) node[below] {$x$};
  \draw[-latex] (0,{(#4)-0.5}) -- (0,{(#5)+0.5}) node[left] {$y$};
  
  % Origin
  \node at (0,0) [below left] {$O$};
}

\makeatother
